% Отчёт по лабораторной работе — кубическая аппроксимация (многочлен Эрмита)
\documentclass[a4paper,12pt]{article}
\usepackage[T2A]{fontenc}
\usepackage[utf8]{inputenc}
\usepackage[russian]{babel}
\usepackage{amsmath,amssymb}
\usepackage{booktabs}
\usepackage{geometry}
\geometry{left=2.5cm,right=2.5cm,top=2.0cm,bottom=2.0cm}

\title{Лабораторная работа — Метод кубической аппроксимации (многочлен Эрмита)}
\date{}

\begin{document}
\maketitle

\section*{Постановка задачи}
Минимизация функции
\[ f(x)=\frac{1}{x}+e^{x}, \qquad x\in[0.5,\;1.5].\]
Используется метод кубической аппроксимации (кубический интерполяционный многочлен Эрмита) на текущем отрезке $[a,b]$. Критерий останова: по условию из методички процесс продолжается, пока
\[\delta = b-a \ge \varepsilon_1,\qquad |f'(x_m)| \ge \varepsilon_2,\]
где $\varepsilon_1=\varepsilon_2=10^{-4}$.

\section*{Короткое описание метода}
На каждом шаге, для текущего отрезка $[x_0,x_1]$ строим кубический Эрмитов многочлен $H(x)$ такой, чтобы
\[ H(x_0)=f(x_0),\quad H'(x_0)=f'(x_0),\quad H(x_1)=f(x_1),\quad H'(x_1)=f'(x_1). \]

Его можно записать в виде
\[ H(x)=\alpha_3 x^3 + \alpha_2 x^2 + \alpha_1 x + \alpha_0. \]

Коэффициенты (обобщённо) вычисляются из системы интерполяции. Для поиска стационарных точек решаем уравнение
\[ H'(x)=3\alpha_3 x^2 + 2\alpha_2 x + \alpha_1 = 0. \]
Отсюда
\begin{equation}\label{eq:xm}
x_m = \frac{-\alpha_2 \pm \sqrt{\alpha_2^2 - 3\alpha_1\alpha_3}}{3\alpha_3}.
\end{equation}

Из полученных корней выбираем тот, который лежит внутри $[x_0,x_1]$ (если таковой есть); вычисляем $f(x_m)$ и $f'(x_m)$. Затем обновляем границы отрезка по правилу:
\begin{align*}
\text{если } f'(x_m)<0 &\quad\Rightarrow\quad x_0 := x_m,\\
\text{если } f'(x_m)>0 &\quad\Rightarrow\quad x_1 := x_m,\\
\text{если } f'(x_m)=0 &\quad\Rightarrow\quad \text{найден стационарный минимум.}
\end{align*}

Алгоритм повторяется, пока одновременно выполняются условия на $\delta$ и $|f'(x_m)|$.

\subsection*{Формулы}
Рассмотрим удобное для вывода формул представление кубического интерполяционного многочлена Эрмита в базисе, связанном с узлами $x_1,x_2$:
\begin{equation}\label{eq:P3_newton}
P_3(x)=a_0 + a_1(x-x_1)+a_2(x-x_1)(x-x_2)+a_3(x-x_1)^2(x-x_2).
\end{equation}
Здесь положим
\[y_1=f(x_1),\qquad y_2=f(x_2),\qquad y_{11}=f'(x_1),\qquad y_{12}=f'(x_2),\qquad h=x_2-x_1.\]
Коэффициенты $a_0,a_1,a_2,a_3$ однозначно выражаются через $y_1,y_2,y_{11},y_{12}$. Из необходимого условия $P_3'(x)=0$, подставив коэффициенты, получаем квадратичное уравнение для смещённой переменной $t=x-x_1$ вида
Для явности выпишем прямо формулы для коэффициентов в терминах значений и производных в концах отрезка. Обозначим $h=x_2-x_1$. Тогда
\begin{align}
a_0 &= y_1, \\
a_1 &= \dfrac{y_2-y_1}{h}, \\
a_2 &= \dfrac{y_2-y_1}{h^2} - \dfrac{y_{11}}{h}, \\
a_3 &= \dfrac{y_{11}+y_{12} - 2\dfrac{y_2-y_1}{h}}{h^2}.
\end{align}

Эти выражения получаются из системы условий интерполяции при подстановке разложения (\ref{eq:P3_newton}) в условия на значения и производные в точках $x_1,x_2$. Теперь, из необходимого условия $P_3'(x)=0$, подставив коэффициенты, получаем квадратичное уравнение для смещённой переменной $t=x-x_1$ вида
\begin{equation}\label{eq:quad_t}
3a_3 t^2 - 2\frac{y_{11}+y_{12}-3\dfrac{y_2-y_1}{h}}{h}\;t + y_{11}=0.
\end{equation}
Решение этого уравнения удобно представить в нормированной форме
\begin{equation}\label{eq:x_mu}
x = x_1 + \mu\,h,
\end{equation}
где дробь $\mu$ выражается через промежуточные величины
\begin{align}
z &= y_{11}+y_{12}-3\dfrac{y_2-y_1}{h},\\
\omega &= \sqrt{z^2 - y_{11}\,y_{12}}.
\end{align}
Тогда одно из компактных представлений корня даётся формулой:
\begin{equation}\label{eq:mu_formula}
\mu = \frac{\omega + z - y_{11}}{2\omega - y_{11} + y_{12}}.
\end{equation}
Подставляя в (\ref{eq:x_mu}) получаем явную формулу для рассматриваемой стационарной точки
\[ x = x_1 + \frac{\omega + z - y_{11}}{2\omega - y_{11} + y_{12}}\, (x_2-x_1). \]
Это представление эквивалентно формуле (\ref{eq:xm}) и удобно тем, что выражение в правой части записано в масштабе отрезка $[x_1,x_2]$ (коэффициент $\mu$ — безразмерная величина). При реализации следует контролировать подкоренное выражение $z^2-y_{11}y_{12}$ (оно должно быть неотрицательным для вещественного корня) и выбирать знак корня, дающий точку внутри $[x_1,x_2]$.

\section*{Аналитические расчёты — первые три итерации}
Ниже приведены результаты первых трёх итераций алгоритма, полученные при запуске программы:

\begin{table}[h]
\centering
% Уменьшаем таблицу: уменьшаем шрифт и межстолбцовый отступ
{\scriptsize
\setlength{\tabcolsep}{4pt}
\begin{tabular}{r r r r r r r r r r}
\toprule
 & $a$ & $b$ & $\Delta$ & $x_m$ & $f(x_m)$ & $f'(x_m)$ & $\alpha_3$ & $\alpha_2$ & $\alpha_1$ \\
\midrule
1 & 0.50000000 & 1.50000000 & 1.000e+00 & 0.75184583 & 3.45097116 & 3.519e-01 & -1.313e+00 & 7.134e+00 & -8.500e+00 \\
2 & 0.50000000 & 0.75184583 & 2.518e-01 & 0.69925214 & 3.44234657 & -3.294e-02 & -4.290e-01 & 1.146e+00 & -9.734e-01 \\
3 & 0.69925214 & 0.75184583 & 5.259e-02 & 0.70347954 & 3.44227730 & 9.407e-05 & -9.055e-03 & 2.983e-02 & -2.852e-02 \\
\bottomrule
\end{tabular}
}
\label{tab:iters}
\end{table}

В таблице также присутствует дискриминант квадратичного уравнения производной $H'(x)=0$; для полноты приведём его (значения взяты из вычислительного отчёта):
\begin{itemize}
  \item итерация 1: $\mathrm{D_1}=1.741 \cdot 10^{1}$;
  \item итерация 2: $\mathrm{D_2}=6.052 \cdot 10^{-2}$;
  \item итерация 3: $\mathrm{D_3}=1.149 \cdot 10^{-4}$.
\end{itemize}

\paragraph{Краткий показательный расчёт для 1-й итерации.}
\begin{enumerate}
  \item Взяли отрезок $[a,b]=[0.5,1.5]$. Вычислили $y_0=f(0.5),\;y_1=f(1.5),\;y_0'=f'(0.5),\;y_1'=f'(1.5)$ (значения используются в коде при построении Эрмита).
  \item Построили $H(x)$ — кубический интерполянт Эрмита по четырём условиям на концах.
  \item Решили квадратное уравнение для $H'(x)=0$ и использовали формулу (\ref{eq:xm}). Получили $x_m=0.75184583$. Вычислили $f(x_m)=3.45097116$ и $f'(x_m)=3.519\times10^{-1}$.
  \item Поскольку $f'(x_m)>0$, обновили правую границу: $b:=x_m$ (см. действие в таблице). Новый отрезок для 2-й итерации: $[0.5,\;0.75184583]$.
\end{enumerate}

\paragraph{Краткие итоги первых трёх итераций.}
После трёх итераций получено приближение
\[ x^* \approx 0.703479536307, \qquad f(x^*) \approx 3.44227729506, \qquad f'(x^*) \approx 9.4072\times10^{-5}.\]
Критерий останова выполнен (одно из условий нарушено — малое значение производной в узле). Алгоритм сошёлся за 3 итерации.

\section*{Выводы}
Реализованный алгоритм соответствует схеме из методички: на каждом шаге строился кубический интерполяционный многочлен Эрмита на концах текущего отрезка, находилась стационарная точка многочлена и по знаку производной в этой точке обновлялись границы отрезка. Для функции $f(x)=1/x+e^{x}$ алгоритм быстро сходится — в рассматриваемом запуске за 3 итерации.
\end{document}